\documentclass[11pt]{article}
\usepackage[margin=1in]{geometry}
\usepackage{amsmath,amssymb,amsthm}
\usepackage{booktabs,longtable,array}
\usepackage{hyperref}
\usepackage{xcolor}

\newtheorem{theorem}{Theorem}[section]
\newtheorem{lemma}[theorem]{Lemma}
\newtheorem{proposition}[theorem]{Proposition}
\newtheorem{corollary}[theorem]{Corollary}
\theoremstyle{definition}
\newtheorem{definition}[theorem]{Definition}
\theoremstyle{remark}
\newtheorem{remark}[theorem]{Remark}

\newcommand{\F}{\mathcal{F}}
\newcommand{\Aut}{\mathrm{Aut}}
\newcommand{\Out}{\mathrm{Out}}
\newcommand{\Inn}{\mathrm{Inn}}
\newcommand{\Syl}{\mathrm{Syl}}
\newcommand{\GL}{\mathrm{GL}}
\newcommand{\SL}{\mathrm{SL}}
\newcommand{\PSL}{\mathrm{PSL}}
\newcommand{\PSU}{\mathrm{PSU}}
\newcommand{\PSp}{\mathrm{PSp}}
\newcommand{\Alt}{\mathrm{Alt}}
\newcommand{\Sym}{\mathrm{Sym}}
\newcommand{\Z}{\mathrm{Z}}
\newcommand{\Phifr}{\Phi}
\newcommand{\gen}[1]{\langle #1 \rangle}

\title{Saturated fusion systems on groups of order $3^8$\\[1ex]
\large a running report}
\author{}
\date{\today}

\begin{document}
\maketitle

\begin{abstract}
We search for the saturated fusion systems $\F$ on the groups $S$ of order
$3^8$ with $O_3(\F)=1$ and $O^3(\F)=\F$, in the manner of
Parker--Semeraro \cite{PS} but with the fast proto-essential search of
Gautam \cite{Gautam} and the fusion-graph-free saturation test of the
GAP package \texttt{FusionSystems} that accompanies this report.  There are
$1\,396\,077$ groups of order $3^8$ (Vaughan-Lee's database, available in
GAP through the package \texttt{sglppow}).  Section~\ref{sec:reduction}
reduces them to a short list by necessary conditions, Section~\ref{sec:groups}
identifies the surviving groups, Section~\ref{sec:lookup} contains the
look-up tables of subgroups of $\GL_n(3)$ with a strongly $3$-embedded
subgroup, and Section~\ref{sec:results} records the fusion systems found
group by group.  This document is updated as the computation proceeds; the
current state is summarised in Section~\ref{sec:status}.
\end{abstract}

\tableofcontents

\section{Status}\label{sec:status}
\textcolor{red}{[updated as the work proceeds]}
\begin{itemize}
\item Stage 1a (cheap reduction of the $1\,396\,077$ groups): \emph{running}.
\item Stage 1b (automorphism group tests on the survivors): not started.
\item Look-up tables for $\GL_n(3)$: complete for $n\le 4$ (Section~\ref{sec:lookup}).
\item Stage 2 (fusion systems group by group): not started.
\end{itemize}

\section{Conventions and the strategy}\label{sec:strategy}

Throughout $p=3$, $S$ is a group of order $3^8$ and $\F$ is a saturated
fusion system on $S$ with $O_3(\F)=1$ and $O^3(\F)=\F$.  We write
$\Aut_S(E)=N_S(E)/C_S(E)$ and $\Out_S(E)=\Aut_S(E)\Inn(E)/\Inn(E)$; for an
$S$-centric subgroup $E$ (that is $C_S(E)\le E$) we have
$\Out_S(E)\cong N_S(E)/E$.

By the Alperin--Goldschmidt fusion theorem $\F$ is generated by $\Aut_\F(S)$
and the groups $\Aut_\F(E)$ for the $\F$-essential subgroups $E$ (the
$\F$-centric, fully normalised subgroups with $\Out_\F(E)$ having a strongly
$3$-embedded subgroup).  Following \cite{PS} we therefore proceed as
follows.
\begin{enumerate}
\item Determine the subgroups $E\le S$ which can be essential in some
  saturated fusion system on $S$ (the \emph{proto-essential} subgroups), up
  to $\Aut(S)$-conjugacy.  Here we use Gautam's theorem (Theorem~\ref{thm:gautam}
  below) instead of the subgroup lattice of $S$.
\item Determine the Borel subgroups, i.e.\ the $3'$-subgroups $K\le\Aut(S)$
  up to conjugacy, so that $\Aut_\F(S)=\Inn(S)K$ (recall that $\Inn(S)$ is a
  normal Sylow $3$-subgroup of $\Aut_\F(S)$).
\item For each Borel subgroup and each proto-essential subgroup $E$ determine
  the candidates for $\Aut_\F(E)$: the subgroups of $\Aut(E)$ containing
  $\Inn(E)$ with $\Aut_S(E)$ as a Sylow $3$-subgroup, with $\Out$ having a
  strongly $3$-embedded subgroup and compatible with the Borel subgroup.
\item For each choice of a set of essential subgroups (up to the action of
  $N_{\Aut(S)}(\Aut_\F(S))$) and automizers, build the fusion system,
  discard it unless $O^3(\F)=\F$, $O_3(\F)=1$ and $\F$ is saturated, and
  finally remove isomorphic copies.
\end{enumerate}
Steps 2--4 are only carried out for the groups surviving the reduction of
Section~\ref{sec:reduction}.  All computations are done with the GAP package
\texttt{FusionSystems}; the scripts for this report are in the directory
\texttt{order38/scripts} of the repository.

\begin{theorem}[Gautam {\cite[Theorem 3.6]{Gautam}}]\label{thm:gautam}
Let $S=L_1>L_2>\dots>L_m=1$ be a central series of $S$ with all factors
of order $p$ which refines the lower central series.  Every essential
subgroup of a saturated fusion system on $S$ is $S$-conjugate to a subgroup
$C_S(xA/A)=\{s\in S : [s,x]\in A\}$ with $A=L_i\le [S,S]$ and $xA$ an
element of order $p$ of $S/A$.
\end{theorem}

\section{The reduction}\label{sec:reduction}

\subsection{Necessary conditions}

We use the following necessary conditions on $S$.  Conditions
(R1)--(R3) involve only $S$, its candidate subgroups $E$ and (for the few
candidates passing the structural tests) $\Aut(E)$; they are applied to all
$1\,396\,077$ groups (Stage 1a).  Conditions (R4)--(R6) need $\Aut(S)$ and
are applied to the survivors (Stage 1b).

\begin{lemma}[{\cite[Lemma 7.1]{PS}}]\label{lem:R1}
{\rm (R1)} If $|S:\Z(S)|\le 9$ then $S$ carries no saturated fusion system
with $O_3(\F)=1$.
\end{lemma}

\begin{lemma}\label{lem:R2}
{\rm (R2)} Let $E$ be an essential subgroup of a saturated fusion system
$\F$ on $S$ and $T=\Out_S(E)\cong N_S(E)/E$.  Then
\begin{enumerate}
\item[(a)] $E$ is $S$-centric and $E<N_S(E)$;
\item[(b)] $T$ is a Sylow $3$-subgroup of a group with a strongly
  $3$-embedded subgroup, so $T$ is cyclic, or elementary abelian of rank at
  least $2$, or $T\cong 3^{1+2}_+$, or $T\cong\Syl_3(\PSU_3(9))$ of order
  $3^6$;
\item[(c)] $|E/\Phifr(E)|\ge |T|^2$;
\item[(d)] if $1=D_0<D_1<\dots<D_k=E$ is a chain of characteristic
  subgroups of $E$, then no element of $N_S(E)\setminus E$ centralises every
  factor $D_{i}/D_{i-1}$;
\item[(e)] $\Aut(E)$ is not a $3$-group, $\Aut_S(E)\cap O_3(\Aut(E))=\Inn(E)$,
  and $\Aut(E)$ is insoluble unless $T$ is cyclic;
\item[(f)] if $T$ is not cyclic then $|\Aut(E)|$ is divisible by $5$ for
  $T\cong 3^2$, by $7\cdot13$ for $T\cong3^3$, by $7$ for $T\cong3^{1+2}_+$,
  by $5\cdot41$ for $T\cong3^4$, by $11\cdot61$ for $T\cong3^5$ and by
  $5\cdot73$ for $|T|=3^6$.
\end{enumerate}
\end{lemma}
\begin{proof}
(a) and (b) are part of the definition and the classification of the
groups with a strongly $p$-embedded subgroup (for $p=3$ and $T$ non-cyclic
the possibilities for $O^{3'}(\Out_\F(E)/O_{3'}(\Out_\F(E)))$ are
$\PSL_2(3^n)$, $\PSU_3(3^n)$, ${}^2G_2(3^n)$, $\PSL_3(4)$ and $M_{11}$; the
Ree groups have Sylow $3$-subgroups of order at least $3^9$).  (c) is
\cite[Lemma 2.?]{PS}: $\Out_\F(E)$ acts faithfully on $E/\Phifr(E)$ because
the kernel of $\Aut(E)\to\Aut(E/\Phifr(E))$ is a normal $3$-subgroup and
$O_3(\Out_\F(E))=1$ as $E$ is $\F$-radical.  (d) The automorphisms of $E$
centralising all factors of a chain of characteristic subgroups form a
normal subgroup $K$ of $\Aut(E)$, and $K$ is a $3$-group (a $3'$-automorphism
centralising all factors of a normal series of a $3$-group is trivial).
Hence $K\Inn(E)/\Inn(E)\le O_3(\Out(E))$ meets $\Out_\F(E)$ in a normal
$3$-subgroup, which is trivial, so $K\cap\Aut_S(E)\le\Inn(E)$, which is
(d) as $E$ is $S$-centric.  (e) is the $\F$-radical condition
$O_3(\Out_\F(E))=1$ together with the fact that $T\ne1$ is not normal in
$\Out_\F(E)$ (a group with a strongly $3$-embedded subgroup has more than
one Sylow $3$-subgroup), so that $\Out(E)$ is not a $3$-group; if $T$ is not
cyclic then $O^{3'}(\Out_\F(E)/O_{3'}(\Out_\F(E)))$ is a non-abelian simple
group by the classification quoted in (b), so $\Aut(E)$ is insoluble.
(f) That simple group is a section of $\Out(E)$ and is one of $\PSL_2(9)$,
$\PSL_3(4)$, $M_{11}$ for $T\cong3^2$, $\PSL_2(27)$ for $3^3$,
$\PSU_3(3)$ for $3^{1+2}_+$, $\PSL_2(81)$ for $3^4$, $\PSL_2(243)$ for
$3^5$ and $\PSL_2(729)$ or $\PSU_3(9)$ for $|T|=3^6$; the primes listed
divide the orders of all groups in the respective list.
\end{proof}

In Stage 1a a subgroup $E$ is called a \emph{candidate} if it is
$S$-conjugate to a subgroup $C_S(xA/A)$ of Theorem~\ref{thm:gautam} and
satisfies (a)--(f) of Lemma~\ref{lem:R2}, where in (d) we use the chain
$E>\Phifr(E)>1$ and then the covering pairs of the set of characteristic
subgroups generated (by one round of intersections and products) by
$\Z(E)$, $\Z_2(E)$, $\Phifr(E)$, $E'$, $\gamma_3(E)$, $\mho_1(E)$,
$\Omega_1(\Z(E))$, $\Z(\Phifr(E))$ and $C_E(E')$.  (The stability group of
the set of covering pairs is contained in the stability group of any
maximal chain in the set, so the argument of (d) applies.)  The tests are
applied in this order, so that $\Aut(E)$ (computed with the GAP package
\texttt{autpgrp}) is only needed for the candidates surviving (a)--(d);
the first part of (e) already discards almost all of them, since almost
all groups of order $3^7$ have a $3$-group of automorphisms.

\begin{lemma}\label{lem:R3}
{\rm (R3)} If $W\ne 1$ is a characteristic subgroup of $S$ which is contained
in and characteristic in every candidate, then $W\le O_3(\F)$ for every
saturated fusion system $\F$ on $S$.
\end{lemma}
\begin{proof}
$W$ is normalised by $\Aut_\F(S)$ and by $\Aut_\F(E)$ for every essential
$E$ (an essential subgroup is $S$-conjugate to a candidate $E_0$, and $W=W^s$
is characteristic in $E_0^s$).  By Alperin--Goldschmidt $W$ is normalised
by every morphism of $\F$ defined on it, i.e.\ $W\trianglelefteq\F$.
\end{proof}
In Stage 1a ``characteristic'' is replaced by the sufficient condition of
membership in the sets of visibly characteristic subgroups described above.

\begin{lemma}\label{lem:R456}
Let $\mathcal{E}$ be a set of representatives of the $S$-classes of the
proto-essential subgroups of $S$: the candidates of Stage 1a, where a
candidate whose $\Aut(S)$-orbit contains a subgroup which is not a
candidate is discarded (all the conditions of Lemma~\ref{lem:R2} are
$\Aut(S)$-invariant).
\begin{enumerate}
\item[(R4)] If $\mathcal{E}=\emptyset$ then $O_3(\F)=S$ for every saturated
  $\F$ on $S$.
\item[(R5)] If a nontrivial characteristic subgroup of $S$ is characteristic
  in every member of $\mathcal{E}$, then $O_3(\F)\neq1$.
\item[(R6)] If $S\ne\gen{[S,O^{3}(\Aut(S))],\,E : E\in\mathcal{E}}$ then
  $O^3(\F)<\F$ for every saturated $\F$ on $S$.
\end{enumerate}
\end{lemma}
\begin{proof}
(R4) and (R5) are as before.  (R6): $O^3(\F)=\F$ if and only if the
hyperfocal subgroup $\gen{[P,O^3(\Aut_\F(P))] : P\le S}$ is $S$, and by
Alperin--Goldschmidt it is generated by $[S,O^3(\Aut_\F(S))]$ and the
$[E,O^3(\Aut_\F(E))]\le E$ for $E$ essential; finally $O^3(\Aut_\F(S))\le
O^3(\Aut(S))$.
\end{proof}

\subsection{Statistics}
\textcolor{red}{[to be filled in when Stage 1a finishes]}

\section{The surviving groups}\label{sec:groups}
\textcolor{red}{[to be filled in]}

\subsection{Sylow 3-subgroups of simple groups of order $3^8$}\label{sec:simple}
The finite simple groups $G$ with $|G|_3=3^8$ and non-abelian Sylow
$3$-subgroups are (up to the obvious isomorphisms):
\begin{itemize}
\item $\Alt(18)$, $\Alt(19)$, $\Alt(20)$ (and the symmetric groups), with
  $S\cong (C_3\wr C_3)\times(C_3\wr C_3)$;
\item $\PSp_4(9)$, $S$ a Sylow $3$-subgroup of $\mathrm{Sp}_4(9)$;
\item $\PSL_3(q)$ with $3^4\,\|\,q-1$ (e.g.\ $q=163$) and $\PSU_3(q)$ with
  $3^4\,\|\,q+1$ (e.g.\ $q=647$), with $S\cong\big(\{(x,y,z)\in C_{81}^3 :
  xyz=1\}\rtimes C_3\big)/C_3$ of maximal class;
\item ${}^3D_4(q)$ with $3^3\,\|\,q^2-1$ (e.g.\ $q=53$), where $S$ is an
  extension of $C_{81}\times C_{27}$ by $C_3$ (not constructed here).
\end{itemize}
The groups $\PSL_2(3^8)$, $\PSL_2(q)$ with $3^8\,\|\,q\pm1$ and $\PSp_4(q)$
with $3^4\,\|\,q^2-1$ have abelian Sylow $3$-subgroups, which carry no
fusion system with $O_3(\F)=1$.  No sporadic group has $|G|_3=3^8$.
The first three groups are constructed in \texttt{scripts/known\_groups.g};
their invariants are
\begin{center}
\begin{tabular}{lcccccc}
\toprule
$S$ & $|\Z(S)|$ & class & rank & exponent & $|S'|$ & classes\\
\midrule
$\Syl_3(\Alt(18))$ & 9 & 3 & 4 & 9 & 81 & 289\\
$\Syl_3(\PSp_4(9))$ & 9 & 3 & 4 & 9 & 81 & 161\\
$\Syl_3(\PSL_3(163))$ & 3 & 7 & 2 & 81 & 729 & 737\\
\bottomrule
\end{tabular}
\end{center}

\section{Look-up tables: subgroups of $\GL_n(3)$ with a strongly $3$-embedded subgroup}\label{sec:lookup}

If $E\cong 3^n$ is an elementary abelian essential subgroup then
$\Out_\F(E)=\Aut_\F(E)\le\GL_n(3)$ has Sylow $3$-subgroup $T=\Aut_S(E)$
of order at most $3^{8-n}$ (as $E$ is centric and $N_S(E)>E$, in fact
$n\le 7$) and a strongly $3$-embedded subgroup.  The following tables list,
for each $\GL_n(3)$-class of admissible $T$ (cyclic, elementary abelian or
$3^{1+2}_+$), the subgroups $H$ with $T\in\Syl_3(H)$ and a strongly
$3$-embedded subgroup, up to conjugacy in $N_{\GL_n(3)}(T)$.  For $n\le4$ the
tables are complete.  They are produced by \texttt{scripts/gltables.g} and
stored in machine-readable form in \texttt{data/gl$n$\_table.g}.

The classes of $T$ are distinguished in the tables by $|N_{\GL_n(3)}(T)|$;
when the same $H$ appears several times for one $T$ the copies are not
conjugate under $N_{\GL_n(3)}(T)$ (they differ in the way $H$ acts, e.g.\ in
the determinant or in the composition factors).  The column ``factors''
gives the dimensions of the composition factors of $H$ on the natural
module, so that e.g.\ $1+2$ for $\SL_2(3)\le\GL_3(3)$ means the module has a
trivial submodule with irreducible $2$-dimensional quotient.  For $n\le 4$
the only candidates with a non-cyclic Sylow $3$-subgroup occur for two
classes of $T\cong 3^2$ in $\GL_4(3)$: the extensions of $\Alt(6)\cong
\Omega_4^-(3)$ and of $\SL_2(9)=2.\Alt(6)$ by subgroups of
$\mathrm{Out}(\Alt(6))$ and by scalars; $T\cong C_9$, $3^3$, $3^4$ and
$3^{1+2}_+$ admit no candidates in $\GL_4(3)$.  (The Ree groups and
$\PSU_3(9)$ need $n\ge 6$, $\PSL_3(4)$ and $M_{11}$ need $n\ge 5$.)

% automatically generated by gltables.g
\begin{longtable}{llrlrrll}
\caption{Subgroups $H$ of $\GL_2(3)$ with a strongly $3$-embedded subgroup, by Sylow $3$-subgroup $T$ (one $H$ per $N_{\GL_2(3)}(T)$-class; $|N_H(T)|$, $|O_{3'}(H)|$, whether $H$ is soluble, the dimensions of the composition factors of the natural module).}\label{tab:gl2}\\
\toprule
$T$ & $|N(T)|$ & $|H|$ & $H$ & $|N_H(T)|$ & $|O_{3'}(H)|$ & sol. & factors\\
\midrule
\endfirsthead
\toprule
$T$ & $|N(T)|$ & $|H|$ & $H$ & $|N_H(T)|$ & $|O_{3'}(H)|$ & sol. & factors\\
\midrule
\endhead
\bottomrule
\endfoot
$C3$ & 12 & 24 & \texttt{SL(2,3)} & 6 & 8 & true & 2\\
$C3$ & 12 & 48 & \texttt{GL(2,3)} & 12 & 8 & true & 2\\
\end{longtable}

% automatically generated by gltables.g
\begin{longtable}{llrlrrll}
\caption{Subgroups $H$ of $\GL_3(3)$ with a strongly $3$-embedded subgroup, by Sylow $3$-subgroup $T$ (one $H$ per $N_{\GL_3(3)}(T)$-class; $|N_H(T)|$, $|O_{3'}(H)|$, whether $H$ is soluble, the dimensions of the composition factors of the natural module).}\label{tab:gl3}\\
\toprule
$T$ & $|N(T)|$ & $|H|$ & $H$ & $|N_H(T)|$ & $|O_{3'}(H)|$ & sol. & factors\\
\midrule
\endfirsthead
\toprule
$T$ & $|N(T)|$ & $|H|$ & $H$ & $|N_H(T)|$ & $|O_{3'}(H)|$ & sol. & factors\\
\midrule
\endhead
\bottomrule
\endfoot
$C3$ & 216 & 24 & \texttt{SL(2,3)} & 6 & 8 & true & 1+2\\
$C3$ & 216 & 48 & \texttt{GL(2,3)} & 12 & 8 & true & 2+1\\
$C3$ & 216 & 48 & \texttt{C2 x SL(2,3)} & 12 & 16 & true & 2+1\\
$C3$ & 216 & 48 & \texttt{GL(2,3)} & 12 & 8 & true & 2+1\\
$C3$ & 216 & 96 & \texttt{C2 x GL(2,3)} & 24 & 16 & true & 1+2\\
$C3$ & 36 & 12 & \texttt{A4} & 3 & 4 & true & 3\\
$C3$ & 36 & 24 & \texttt{C2 x A4} & 6 & 8 & true & 3\\
$C3$ & 36 & 24 & \texttt{S4} & 6 & 4 & true & 3\\
$C3$ & 36 & 24 & \texttt{S4} & 6 & 4 & true & 3\\
$C3$ & 36 & 39 & \texttt{C13 : C3} & 3 & 13 & true & 3\\
$C3$ & 36 & 48 & \texttt{C2 x S4} & 12 & 8 & true & 3\\
$C3$ & 36 & 78 & \texttt{C2 x (C13 : C3)} & 6 & 26 & true & 3\\
$3^2$ & 864 & \multicolumn{6}{l}{none}\\
$3^2$ & 864 & \multicolumn{6}{l}{none}\\
$3^2$ & 108 & \multicolumn{6}{l}{none}\\
$3^{1+2}_+$ & 216 & \multicolumn{6}{l}{none}\\
\end{longtable}

% automatically generated by gltables.g
\begin{longtable}{llrlrrll}
\caption{Subgroups $H$ of $\GL_4(3)$ with a strongly $3$-embedded subgroup, by Sylow $3$-subgroup $T$ (one $H$ per $N_{\GL_4(3)}(T)$-class; $|N_H(T)|$, $|O_{3'}(H)|$, whether $H$ is soluble, the dimensions of the composition factors of the natural module).}\label{tab:gl4}\\
\toprule
$T$ & $|N(T)|$ & $|H|$ & $H$ & $|N_H(T)|$ & $|O_{3'}(H)|$ & sol. & factors\\
\midrule
\endfirsthead
\toprule
$T$ & $|N(T)|$ & $|H|$ & $H$ & $|N_H(T)|$ & $|O_{3'}(H)|$ & sol. & factors\\
\midrule
\endhead
\bottomrule
\endfoot
$C3$ & 46656 & 24 & \texttt{SL(2,3)} & 6 & 8 & true & 2+1+1\\
$C3$ & 46656 & 48 & \texttt{GL(2,3)} & 12 & 8 & true & 1+1+2\\
$C3$ & 46656 & 48 & \texttt{GL(2,3)} & 12 & 8 & true & 1+1+2\\
$C3$ & 46656 & 48 & \texttt{C2 x SL(2,3)} & 12 & 16 & true & 1+2+1\\
$C3$ & 46656 & 48 & \texttt{C2 x SL(2,3)} & 12 & 16 & true & 2+1+1\\
$C3$ & 46656 & 48 & \texttt{GL(2,3)} & 12 & 8 & true & 1+1+2\\
$C3$ & 46656 & 96 & \texttt{C2 x GL(2,3)} & 24 & 16 & true & 2+1+1\\
$C3$ & 46656 & 96 & \texttt{C4 x SL(2,3)} & 24 & 32 & true & 2+2\\
$C3$ & 46656 & 96 & \texttt{Q8 : (C3 : C4)} & 24 & 16 & true & 2+2\\
$C3$ & 46656 & 96 & \texttt{C2 x C2 x SL(2,3)} & 24 & 32 & true & 1+1+2\\
$C3$ & 46656 & 96 & \texttt{C2 x GL(2,3)} & 24 & 16 & true & 1+1+2\\
$C3$ & 46656 & 96 & \texttt{C2 x GL(2,3)} & 24 & 16 & true & 2+1+1\\
$C3$ & 46656 & 96 & \texttt{C2 x GL(2,3)} & 24 & 16 & true & 2+1+1\\
$C3$ & 46656 & 192 & \texttt{SL(2,3) x Q8} & 48 & 64 & true & 2+2\\
$C3$ & 46656 & 192 & \texttt{Q8 : D24} & 48 & 32 & true & 2+2\\
$C3$ & 46656 & 192 & \texttt{C8 x SL(2,3)} & 48 & 64 & true & 2+2\\
$C3$ & 46656 & 192 & \texttt{Q8 : (C3 : C8)} & 48 & 32 & true & 2+2\\
$C3$ & 46656 & 192 & \texttt{Q8 : (C3 : Q8)} & 48 & 32 & true & 2+2\\
$C3$ & 46656 & 192 & \texttt{C4 x GL(2,3)} & 48 & 32 & true & 2+2\\
$C3$ & 46656 & 192 & \texttt{SL(2,3) x D8} & 48 & 64 & true & 2+2\\
$C3$ & 46656 & 192 & \texttt{Q8 : ((C6 x C2) : C2)} & 48 & 32 & true & 2+2\\
$C3$ & 46656 & 192 & \texttt{C2 x C2 x GL(2,3)} & 48 & 32 & true & 1+1+2\\
$C3$ & 46656 & 384 & \texttt{GL(2,3) x Q8} & 96 & 64 & true & 2+2\\
$C3$ & 46656 & 384 & \texttt{SL(2,3) x QD16} & 96 & 128 & true & 2+2\\
$C3$ & 46656 & 384 & \texttt{GL(2,3) x D8} & 96 & 64 & true & 2+2\\
$C3$ & 46656 & 384 & \texttt{Q8 : ((C3 : Q8) : C2)} & 96 & 64 & true & 2+2\\
$C3$ & 46656 & 384 & \texttt{C8 : GL(2,3)} & 96 & 64 & true & 2+2\\
$C3$ & 46656 & 384 & \texttt{C8 x GL(2,3)} & 96 & 64 & true & 2+2\\
$C3$ & 46656 & 384 & \texttt{Q8 : GL(2,3)} & 96 & 64 & true & 2+2\\
$C3$ & 46656 & 768 & \texttt{GL(2,3) x QD16} & 192 & 128 & true & 2+2\\
$C3$ & 648 & 12 & \texttt{A4} & 3 & 4 & true & 1+3\\
$C3$ & 648 & 24 & \texttt{SL(2,3)} & 6 & 8 & true & 2+2\\
$C3$ & 648 & 24 & \texttt{S4} & 6 & 4 & true & 1+3\\
$C3$ & 648 & 24 & \texttt{S4} & 6 & 4 & true & 3+1\\
$C3$ & 648 & 24 & \texttt{C2 x A4} & 6 & 8 & true & 3+1\\
$C3$ & 648 & 24 & \texttt{C2 x A4} & 6 & 8 & true & 3+1\\
$C3$ & 648 & 24 & \texttt{S4} & 6 & 4 & true & 1+3\\
$C3$ & 648 & 24 & \texttt{C2 x A4} & 6 & 8 & true & 3+1\\
$C3$ & 648 & 24 & \texttt{S4} & 6 & 4 & true & 3+1\\
$C3$ & 648 & 39 & \texttt{C13 : C3} & 3 & 13 & true & 3+1\\
$C3$ & 648 & 48 & \texttt{C2 x S4} & 12 & 8 & true & 1+3\\
$C3$ & 648 & 48 & \texttt{C2 x C2 x A4} & 12 & 16 & true & 1+3\\
$C3$ & 648 & 48 & \texttt{C2 x S4} & 12 & 8 & true & 3+1\\
$C3$ & 648 & 48 & \texttt{C2 x S4} & 12 & 8 & true & 1+3\\
$C3$ & 648 & 48 & \texttt{C2 x S4} & 12 & 8 & true & 3+1\\
$C3$ & 648 & 48 & \texttt{C2 x S4} & 12 & 8 & true & 1+3\\
$C3$ & 648 & 48 & \texttt{C2 x S4} & 12 & 8 & true & 3+1\\
$C3$ & 648 & 48 & \texttt{GL(2,3)} & 12 & 8 & true & 2+2\\
$C3$ & 648 & 60 & \texttt{A5} & 6 & 1 & false & 4\\
$C3$ & 648 & 78 & \texttt{C2 x (C13 : C3)} & 6 & 26 & true & 1+3\\
$C3$ & 648 & 78 & \texttt{C2 x (C13 : C3)} & 6 & 26 & true & 1+3\\
$C3$ & 648 & 78 & \texttt{C2 x (C13 : C3)} & 6 & 26 & true & 1+3\\
$C3$ & 648 & 96 & \texttt{((C2 x C2 x C2) : (C2 x C2)) : C3} & 6 & 32 & true & 4\\
$C3$ & 648 & 96 & \texttt{C2 x C2 x S4} & 24 & 16 & true & 3+1\\
$C3$ & 648 & 120 & \texttt{S5} & 12 & 1 & false & 4\\
$C3$ & 648 & 120 & \texttt{C2 x A5} & 12 & 2 & false & 4\\
$C3$ & 648 & 120 & \texttt{S5} & 12 & 1 & false & 4\\
$C3$ & 648 & 156 & \texttt{C2 x C2 x (C13 : C3)} & 12 & 52 & true & 3+1\\
$C3$ & 648 & 192 & \texttt{(C2 x C2 x C2) : S4} & 12 & 32 & true & 4\\
$C3$ & 648 & 192 & \texttt{Q8 : S4} & 12 & 32 & true & 4\\
$C3$ & 648 & 192 & \texttt{((C4 x C2) : D8) : C3} & 12 & 64 & true & 4\\
$C3$ & 648 & 240 & \texttt{C2 x S5} & 24 & 2 & false & 4\\
$C3$ & 648 & 384 & \texttt{(C2 x C2 x C2 x C2) : S4} & 24 & 64 & true & 4\\
$C3$ & 7776 & 24 & \texttt{SL(2,3)} & 6 & 8 & true & 2+2\\
$C3$ & 7776 & 24 & \texttt{SL(2,3)} & 6 & 8 & true & 1+2+1\\
$C3$ & 7776 & 48 & \texttt{GL(2,3)} & 12 & 8 & true & 2+2\\
$C3$ & 7776 & 48 & \texttt{C2 . S4 = SL(2,3) . C2} & 12 & 8 & true & 4\\
$C3$ & 7776 & 48 & \texttt{((C4 x C2) : C2) : C3} & 12 & 16 & true & 4\\
$C3$ & 7776 & 48 & \texttt{C2 x SL(2,3)} & 12 & 16 & true & 1+2+1\\
$C3$ & 7776 & 48 & \texttt{C2 x SL(2,3)} & 12 & 16 & true & 2+2\\
$C3$ & 7776 & 48 & \texttt{GL(2,3)} & 12 & 8 & true & 2+2\\
$C3$ & 7776 & 48 & \texttt{GL(2,3)} & 12 & 8 & true & 2+1+1\\
$C3$ & 7776 & 48 & \texttt{GL(2,3)} & 12 & 8 & true & 1+2+1\\
$C3$ & 7776 & 96 & \texttt{((C2 x C2 x C2) : (C2 x C2)) : C3} & 24 & 32 & true & 4\\
$C3$ & 7776 & 96 & \texttt{((C8 x C2) : C2) : C3} & 24 & 32 & true & 4\\
$C3$ & 7776 & 96 & \texttt{SL(2,3) : C4} & 24 & 16 & true & 4\\
$C3$ & 7776 & 96 & \texttt{(C2 . S4 = SL(2,3) . C2) : C2} & 24 & 16 & true & 4\\
$C3$ & 7776 & 96 & \texttt{((C2 x Q8) : C2) : C3} & 24 & 32 & true & 4\\
$C3$ & 7776 & 96 & \texttt{((C4 x C2) : C2) : S3} & 24 & 16 & true & 4\\
$C3$ & 7776 & 96 & \texttt{((C4 x C2) : C2) : S3} & 24 & 16 & true & 4\\
$C3$ & 7776 & 96 & \texttt{C2 x GL(2,3)} & 24 & 16 & true & 2+2\\
$C3$ & 7776 & 96 & \texttt{(C2 x Q8) : S3} & 24 & 16 & true & 4\\
$C3$ & 7776 & 96 & \texttt{C2 x GL(2,3)} & 24 & 16 & true & 1+1+2\\
$C3$ & 7776 & 120 & \texttt{SL(2,5)} & 12 & 2 & false & 4\\
$C3$ & 7776 & 192 & \texttt{(Q8 x Q8) : C3} & 12 & 64 & true & 2+2\\
$C3$ & 7776 & 192 & \texttt{((C2 x C2 x C2) : (C2 x C2)) : S3} & 48 & 32 & true & 4\\
$C3$ & 7776 & 192 & \texttt{((C4 x C2) : C2) : D12} & 48 & 32 & true & 4\\
$C3$ & 7776 & 192 & \texttt{((C2 . S4 = SL(2,3) . C2) : C2) : C2} & 48 & 32 & true & 4\\
$C3$ & 7776 & 192 & \texttt{((C2 . S4 = SL(2,3) . C2) : C2) : C2} & 48 & 32 & true & 4\\
$C3$ & 7776 & 192 & \texttt{(C2 . S4 = SL(2,3) . C2) : C4} & 48 & 32 & true & 4\\
$C3$ & 7776 & 192 & \texttt{SL(2,3) : D8} & 48 & 32 & true & 4\\
$C3$ & 7776 & 192 & \texttt{((C8 : C2) : (C2 x C2)) : C3} & 48 & 64 & true & 4\\
$C3$ & 7776 & 240 & \texttt{SL(2,5) : C2} & 24 & 4 & false & 4\\
$C3$ & 7776 & 240 & \texttt{C2 . S5 = SL(2,5) . C2} & 24 & 2 & false & 4\\
$C3$ & 7776 & 240 & \texttt{SL(2,5) : C2} & 24 & 2 & false & 4\\
$C3$ & 7776 & 384 & \texttt{(Q8 x Q8) : S3} & 24 & 64 & true & 4\\
$C3$ & 7776 & 384 & \texttt{Q8 : GL(2,3)} & 24 & 64 & true & 2+2\\
$C3$ & 7776 & 384 & \texttt{((C2 . ((C2 x C2 x C2) : C4) = (C4 x C2) . (C4 x C2)) : C2) : C3} & 24 & 128 & true & 4\\
$C3$ & 7776 & 384 & \texttt{((C8 x C2) : C2) : D12} & 96 & 64 & true & 4\\
$C3$ & 7776 & 480 & \texttt{C8 . A5 = SL(2,5) . C4} & 48 & 8 & false & 4\\
$C3$ & 7776 & 480 & \texttt{SL(2,5) : (C2 x C2)} & 48 & 4 & false & 4\\
$C3$ & 7776 & 480 & \texttt{(C2 . S5 = SL(2,5) . C2) : C2} & 48 & 4 & false & 4\\
$C3$ & 7776 & 768 & \texttt{(Q8 x Q8) : D12} & 48 & 128 & true & 4\\
$C3$ & 7776 & 960 & \texttt{(C8 . A5 = SL(2,5) . C4) : C2} & 96 & 8 & false & 4\\
$C3$ & 7776 & 1920 & \texttt{C2 . ((C2 x C2 x C2 x C2) : A5)} & 48 & 32 & false & 4\\
$C3$ & 7776 & 3840 & \texttt{(C2 . ((C2 x C2 x C2 x C2) : A5)) : C2} & 96 & 32 & false & 4\\
$3^2$ & 432 & \multicolumn{6}{l}{none}\\
$3^2$ & 46656 & \multicolumn{6}{l}{none}\\
$3^2$ & 648 & \multicolumn{6}{l}{none}\\
$3^2$ & 648 & \multicolumn{6}{l}{none}\\
$3^2$ & 1944 & \multicolumn{6}{l}{none}\\
$3^2$ & 46656 & \multicolumn{6}{l}{none}\\
$3^2$ & 5832 & \multicolumn{6}{l}{none}\\
$C9$ & 324 & \multicolumn{6}{l}{none}\\
$3^2$ & 10368 & 720 & \texttt{SL(2,9)} & 72 & 2 & false & 4\\
$3^2$ & 10368 & 1440 & \texttt{SL(2,9) : C2} & 144 & 4 & false & 4\\
$3^2$ & 10368 & 1440 & \texttt{SL(2,9) : C2} & 144 & 2 & false & 4\\
$3^2$ & 10368 & 2880 & \texttt{SL(2,9) : C4} & 288 & 8 & false & 4\\
$3^2$ & 10368 & 2880 & \texttt{SL(2,9) : (C2 x C2)} & 288 & 4 & false & 4\\
$3^2$ & 10368 & 2880 & \texttt{SL(2,9) : C4} & 288 & 4 & false & 4\\
$3^2$ & 10368 & 5760 & \texttt{SL(2,9) : D8} & 576 & 8 & false & 4\\
$3^2$ & 10368 & 5760 & \texttt{SL(2,9) : Q8} & 576 & 8 & false & 4\\
$3^2$ & 10368 & 5760 & \texttt{GL(2,9)} & 576 & 8 & false & 4\\
$3^2$ & 10368 & 11520 & \texttt{SL(2,9) : QD16} & 1152 & 8 & false & 4\\
$3^2$ & 2592 & \multicolumn{6}{l}{none}\\
$3^2$ & 864 & 360 & \texttt{A6} & 36 & 1 & false & 4\\
$3^2$ & 864 & 720 & \texttt{C2 x A6} & 72 & 2 & false & 4\\
$3^2$ & 864 & 720 & \texttt{A6 . C2} & 72 & 1 & false & 4\\
$3^2$ & 864 & 720 & \texttt{S6} & 72 & 1 & false & 4\\
$3^2$ & 864 & 1440 & \texttt{A6 : C4} & 144 & 2 & false & 4\\
$3^2$ & 864 & 1440 & \texttt{C2 x S6} & 144 & 2 & false & 4\\
$3^2$ & 864 & 1440 & \texttt{C2 x (A6 . C2)} & 144 & 2 & false & 4\\
$3^2$ & 864 & 2880 & \texttt{(C2 x (A6 . C2)) : C2} & 288 & 2 & false & 4\\
$3^3$ & 5832 & \multicolumn{6}{l}{none}\\
$3^{1+2}_+$ & 3888 & \multicolumn{6}{l}{none}\\
$3^{1+2}_+$ & 2916 & \multicolumn{6}{l}{none}\\
$3^3$ & 7776 & \multicolumn{6}{l}{none}\\
$3^{1+2}_+$ & 1944 & \multicolumn{6}{l}{none}\\
$3^3$ & 3888 & \multicolumn{6}{l}{none}\\
$3^{1+2}_+$ & 3888 & \multicolumn{6}{l}{none}\\
$3^3$ & 7776 & \multicolumn{6}{l}{none}\\
$3^3$ & 606528 & \multicolumn{6}{l}{none}\\
$3^{1+2}_+$ & 23328 & \multicolumn{6}{l}{none}\\
$3^{1+2}_+$ & 972 & \multicolumn{6}{l}{none}\\
$3^3$ & 5832 & \multicolumn{6}{l}{none}\\
$3^{1+2}_+$ & 1944 & \multicolumn{6}{l}{none}\\
$3^3$ & 11664 & \multicolumn{6}{l}{none}\\
$3^{1+2}_+$ & 1944 & \multicolumn{6}{l}{none}\\
$3^{1+2}_+$ & 972 & \multicolumn{6}{l}{none}\\
$3^3$ & 606528 & \multicolumn{6}{l}{none}\\
$3^{1+2}_+$ & 1944 & \multicolumn{6}{l}{none}\\
$3^4$ & 186624 & \multicolumn{6}{l}{none}\\
\end{longtable}



\section{Results group by group}\label{sec:results}
\textcolor{red}{[to be filled in: one subsection per group, in the order of
the computation (groups without large elementary abelian subgroups first),
with a table of the fusion systems found in the style of \cite{PS}: for each
$\F$ the group $\Out_\F(S)$, the essential subgroups with their automizers,
and an identification or the remark ``exotic?'']}

\begin{thebibliography}{9}
\bibitem{PS} C. Parker, J. Semeraro, \emph{Algorithms for fusion systems with
applications to $p$-groups of small order}, Math. Comp. 90 (2021), arXiv:2003.01600.
\bibitem{Gautam} P. Gautam, \emph{Fusion systems on Sylow 3-subgroups of
Fischer and Monster sporadic groups II}, arXiv:2607.24674.
\bibitem{BCGLO} C. Broto, N. Castellana, J. Grodal, R. Levi, B. Oliver,
\emph{Subgroup families controlling $p$-local finite groups}, Proc. LMS 91 (2005).
\end{thebibliography}

\end{document}
